\section{Discusión (Entrega 3)}
\label{sec:discusion}

\subsection{Sobre la fragmentación por fecha frente al patrón de acceso real}

El rediseño de la fragmentación descrito en la Sección~\ref{sec:diseno-esquema} ilustra una tensión
frecuente en el diseño de sistemas distribuidos: la clave de fragmentación que exige el requisito
formal de la asignatura (\texttt{fecha\_apertura}) no coincide con el filtro más frecuente del
dominio real del problema (\texttt{zone}, usado por el control de acceso por rol de técnico y por
la mayoría de los reportes operativos). Se priorizó el cumplimiento literal del requisito sobre la
preferencia de diseño original, documentando explícitamente en el ADR-0003 el costo aceptado
—consultas por zona convertidas en \emph{scatter-gather} sobre las cuatro particiones— en vez de
ocultarlo o de forzar un diseño híbrido no solicitado. En un sistema de producción real, esta
decisión probablemente se resolvería con una fragmentación compuesta (por ejemplo, subparticiones
por zona dentro de cada partición temporal) o con un índice secundario global cubriente sobre
\texttt{zone}; ambas alternativas quedan fuera del alcance de esta entrega pero se identifican como
extensión natural del diseño actual.

\subsection{Sobre la colocalización cliente–ticket entre microservicios}

La limitación más significativa frente a la letra de la guía es la imposibilidad de colocalizar
físicamente \texttt{tickets} con la tabla de clientes, porque esta última es propiedad de
\texttt{auth-service} en una base de datos distinta. Esta tensión —requisito de la guía de
Entrega 3 (pensada para un esquema monolítico o de bases de datos compartidas) contra un principio
central de la arquitectura de microservicios adoptada desde la Entrega 2 (cada servicio es dueño
exclusivo de su propio esquema)— se resolvió documentando la limitación en vez de comprometer el
límite de propiedad de datos entre servicios, que es una de las garantías arquitectónicas más
valiosas de la descomposición en microservicios~\cite{newman2021microservices}: permite evolucionar
o escalar \texttt{ticket-service} sin coordinar cambios de esquema con \texttt{auth-service}. La
alternativa de aplicación —resolución de identidad del cliente vía JWT en vez de una clave física
compartida— preserva esta garantía a costa de no poder ejecutar un \emph{join} físico dentro del
motor de base de datos entre ambas tablas.

\subsection{Sobre la instrumentación y la observabilidad}

La combinación de pruebas de integración con Testcontainers y métricas Prometheus verificadas con
\texttt{promtool} (Sección~\ref{sec:cluster-verificacion}) resultó más valiosa de lo previsto
inicialmente para diagnosticar el propio incidente de saturación de memoria descrito en esa
sección: sin las métricas de reintentos y sin comprender el comportamiento esperado del aislamiento
serializable —verificado previamente en las pruebas de integración—, habría sido más difícil
distinguir un problema de recursos del sistema operativo (memoria insuficiente) de un problema de
diseño en la capa de datos. Esto refuerza un principio general de sistemas distribuidos: la
observabilidad no es solo un requisito de cumplimiento de la rúbrica, sino una herramienta de
diagnóstico que se vuelve indispensable precisamente cuando el sistema falla de forma inesperada.

\subsection{Sobre la representatividad del pipeline analítico}

El pipeline de Spark descrito en la Sección~\ref{sec:pipeline-paralelo} se diseñó deliberadamente
para reflejar un caso de uso analítico propio del dominio (reincidencia y \emph{clustering} de
incidencias) en vez de aplicar las cinco transformaciones exigidas de forma genérica sobre datos
sin relación con el negocio. Esto tiene un costo: los resultados de \emph{speedup} del protocolo
experimental (Sección~\ref{sec:protocolo-resultados}) son específicos a un pipeline dominado por
una etapa de \emph{join}/ventana con \emph{shuffle} moderado y una etapa final de aprendizaje
automático relativamente costosa (vectorización TF-IDF + \texttt{KMeans}), y no necesariamente
generalizan a pipelines de Spark dominados por \emph{shuffles} masivos entre tablas de tamaño
comparable. Se documenta esta amenaza a la validez externa explícitamente en la
Sección~\ref{sec:protocolo-resultados} en vez de presentar la conclusión del experimento como
universalmente aplicable a cualquier carga de trabajo de Spark.

\subsection{Alcance no cubierto en esta entrega}

Dos elementos mencionados en la guía de la asignatura —un \emph{API Gateway} explícito como punto
de entrada único al conjunto de microservicios, y la orquestación completa de los cinco
microservicios mediante un único archivo \texttt{docker-compose.yml} raíz— no se implementaron en
esta entrega y quedan identificados como extensión pendiente: actualmente cada microservicio se
levanta y orquesta de forma independiente, sin una capa de enrutamiento y políticas transversales
(limitación de tasa, agregación de respuestas) centralizada.

\paragraph{Resuelto en la Entrega 4.} Ambos pendientes se cerraron en el ciclo siguiente:
\texttt{api-gateway} (Spring Cloud Gateway) es hoy el único punto de entrada para los dos
clientes, y un único \texttt{docker-compose.yml} en la raíz del repositorio orquesta los 17
contenedores del sistema completo, incluyendo un servicio \texttt{db-init} de un solo uso que
automatiza la inicialización de bases de datos antes manual (Sección~\ref{sec:arquitectura-e4}).
